% Generated draft: requires mathematical and source review.
\documentclass[11pt]{article}
\usepackage[margin=25mm]{geometry}
\usepackage{fontspec}
\setmainfont{Latin Modern Roman}
\usepackage{amsmath,amssymb,mathtools,unicode-math}
\setmathfont{Latin Modern Math}
\usepackage{enumitem,needspace,xurl,hyperref}
\setlength{\parindent}{0pt}
\setlength{\parskip}{5pt}
\setlength{\emergencystretch}{3em}
\newcommand{\N}{\mathbb{N}}
\newcommand{\Z}{\mathbb{Z}}
\newcommand{\Q}{\mathbb{Q}}
\newcommand{\R}{\mathbb{R}}
\newcommand{\C}{\mathbb{C}}
\newcommand{\F}{\mathbb{F}}
\newcommand{\A}{\mathbb{A}}
\newcommand{\PP}{\mathbb{P}}
\newcommand{\E}{\mathbb{E}}
\newcommand{\eps}{\varepsilon}
\newcommand{\dd}{\,\mathrm{d}}
\begin{document}
\section*{UnsolvedMath ID: 6700044}
\textbf{Source problem code:} catalogue entry
\subsection{Definitions and mathematical statement}
Let \((X_0,g_0)\) be a compact area-extremal Riemannian manifold in the sense of the source. Give \(X_0\times\mathbb R^m\) the product of \(g_0\) with the Euclidean metric. The stabilization conjecture has two assertions:
\[
 \begin{array}{ll}
 \mathrm{A}:&X_0\times\mathbb R^m\text{ is area-gap-extremal for every Euclidean factor dimension }m;\\
 \mathrm{B}:&X_0\times\mathbb R\text{ is area-extremal}.
 \end{array}
\]
The two predicates must remain distinct. Their definitions are not included, so no precise scalar-curvature threshold or allowed comparison class can be inferred. In standard comparison terminology an area-nonincreasing smooth map has \(\|\Lambda^2\dd f\|\le1\), but the excerpt does not state the full map, degree, spin, boundary, or curvature conventions defining either extremality property.

The source says the claims can be confirmed for its known examples proved by Dirac-operator methods, and sketches part A for the unit sphere \(X_0=S^n\). It considers a complete orientable spin manifold \(X\) and a smooth proper map \(f:X\to S^n\times\mathbb R^m\), using \(\operatorname{Scal}(\mathbb R^m)=0\) and \(K\text{-}\operatorname{Area}(\mathbb R^m)=\infty\). It assumes \(n\) even, takes the positive spinor bundle \(S^+\) on \(S^n\), and a rationally homologically essential unitary bundle \(L_{\varepsilon_0}\) with an \(\varepsilon_0\)-flat connection flat at infinity. The qualifier on \(m\) is garbled in the extraction as ``misand'' and is not reconstructed.

The indicated twist is the pullback of the exterior tensor product of \(S^+\) and \(L_{\varepsilon_0}\). The sketch uses a lower scalar-curvature bound \(n(n-1)+\varepsilon\) at infinity with \(\varepsilon_0\) sufficiently small relative to \(\varepsilon>0\), nonzero degree for \(f\), and a relative index, then invokes a global lower bound to obtain a contradiction by a Llarull-type calculation. These are attributed proof-sketch ingredients, not a complete proof: the required area-control hypothesis, relative/Fredholm setup, meaning of homological essentialness and flatness at infinity, and quantitative error bound are absent. The text's ``strict positivity'' of a Dirac operator also needs the intended square/Weitzenb\"ock interpretation checked. Its wording repeats ``at infinity'' beside ``on all of \(X\)'', so the two lower-bound stages require source confirmation.

The final sentence starts a separate ``Subcomplete Extremality'' topic with an open, not necessarily complete manifold but supplies no complete additional question. The registry reports partial progress in Dirac-based examples while leaving arbitrary area-extremal \(X_0\) open; neither that status nor the sketch is independently verified here.
\subsection{Short English statement}
Should area extremality survive a product with one Euclidean line, and should a weaker area-gap version survive every Euclidean factor? The source gives a Dirac-based sphere sketch, but omits the definitions and several hypotheses needed for a complete stabilization theorem.
\subsection{Sources}
\begin{itemize}
\item Local registry: \url{data/unsolvedmath/problems.json}, permanent UnsolvedMath ID \(6700044\), source code catalogue entry; complete supplied statement, background, and research summary, updated 2026-08-06.
\item Misha Gromov, \emph{101 Questions, Problems and Conjectures around Scalar Curvature} (2017), printed tag [?47], occurrence 46, PDF page 53: \url{https://www.ihes.fr/~gromov/wp-content/uploads/2018/08/101-problemsOct1-2017.pdf}. Supplied normalized PDF-text extraction accessed 2026-07-29; described as incomplete.
\item Llarull's curvature computation and previous Dirac arguments are referenced in the passage without full bibliographic details or the preceding definitions.
\end{itemize}
\end{document}
